\subsection{Protocolo MCP (Model Context Protocol)}

\subsubsection{Origen y motivación}

A medida que los LLMs se integraron en aplicaciones de producción, surgió un
problema recurrente: cada fuente de datos o herramienta externa requería una
implementación personalizada para conectarse con el modelo. Un agente que
necesitaba acceder a una base de datos, un sistema de archivos y una API de
terceros debía implementar tres conectores distintos, cada uno con su propia
lógica de autenticación, serialización y manejo de errores. Esta fragmentación
hacía que la construcción de sistemas verdaderamente conectados fuera difícil de
escalar \citep{anthropic2024mcp}.

En noviembre de 2024, Anthropic introdujo el Model Context Protocol (MCP) como
un estándar abierto para resolver este problema. MCP proporciona un protocolo
uniforme para conectar aplicaciones de inteligencia artificial con los sistemas
donde residen los datos, estableciendo conexiones bidireccionales y seguras
entre fuentes de datos y aplicaciones basadas en LLMs
\citep{anthropic2024mcp}. La analogía oficial del protocolo lo compara con un
puerto USB-C: así como USB-C estandarizó la conexión entre dispositivos
electrónicos, MCP estandariza la conexión entre aplicaciones de IA y sistemas
externos \citep{mcpspec2024}.

\subsubsection{Arquitectura cliente-servidor}

MCP sigue una arquitectura cliente-servidor con tres participantes principales
\citep{mcpspec2024}:

\begin{itemize}
    \item \textbf{MCP Host}: la aplicación de IA que coordina y gestiona uno o
    múltiples clientes MCP. Ejemplos incluyen Claude Desktop, Visual Studio
    Code con Copilot, Cursor y Claude Code.

    \item \textbf{MCP Client}: un componente dentro del host que mantiene una
    conexión dedicada con un servidor MCP específico. El host instancia un
    cliente por cada servidor al que se conecta.

    \item \textbf{MCP Server}: un programa que expone datos, herramientas o
    plantillas de interacción a los clientes MCP. Los servidores pueden
    ejecutarse localmente (en la misma máquina que el host) o de forma remota
    (en un servidor web).
\end{itemize}

Esta separación permite que un mismo host se conecte simultáneamente a
múltiples servidores ---por ejemplo, un servidor de sistema de archivos local,
un servidor de base de datos y un servidor remoto de Sentry--- manteniendo
conexiones independientes y aisladas para cada uno.

\subsubsection{Capas del protocolo}

El protocolo se estructura en dos capas conceptuales \citep{mcpspec2024}:

\paragraph{Capa de datos.} Implementa un protocolo de intercambio basado en
JSON-RPC 2.0 \citep{jsonrpc2010spec} que define la estructura y semántica de
los mensajes. Esta capa incluye:

\begin{itemize}
    \item \textbf{Gestión del ciclo de vida}: manejo de la inicialización de
    conexión, negociación de capacidades y terminación. Durante la
    inicialización, cliente y servidor intercambian sus capacidades soportadas
    (herramientas, recursos, prompts) y versión del protocolo mediante un
    handshake formal.

    \item \textbf{Primitivas de servidor}: los tres tipos fundamentales de
    contexto que un servidor puede exponer:
    \begin{itemize}
        \item \textit{Tools} (herramientas): funciones ejecutables que la
        aplicación de IA puede invocar para realizar acciones ---operaciones
        de archivos, consultas a bases de datos, llamadas a APIs.
        \item \textit{Resources} (recursos): fuentes de datos que proveen
        información contextual ---contenido de archivos, registros de base
        de datos, respuestas de APIs.
        \item \textit{Prompts} (plantillas): templates reutilizables que
        estructuran las interacciones con el modelo.
    \end{itemize}

    \item \textbf{Primitivas de cliente}: capacidades que el cliente expone al
    servidor, incluyendo \textit{sampling} (solicitar completions al LLM del
    host), \textit{elicitation} (solicitar información adicional al usuario) y
    \textit{logging} (enviar mensajes de depuración).

    \item \textbf{Notificaciones}: mensajes en tiempo real que permiten a los
    servidores informar a los clientes sobre cambios sin necesidad de polling
    ---por ejemplo, cuando la lista de herramientas disponibles cambia
    dinámicamente.
\end{itemize}

\paragraph{Capa de transporte.} Define los mecanismos de comunicación entre
clientes y servidores. MCP soporta dos mecanismos de transporte:

\begin{itemize}
    \item \textbf{Stdio}: utiliza flujos de entrada/salida estándar para
    comunicación directa entre procesos locales en la misma máquina. Ofrece
    rendimiento óptimo sin overhead de red y es el mecanismo típico para
    servidores locales.

    \item \textbf{Streamable HTTP}: utiliza HTTP POST para mensajes del cliente
    al servidor con Server-Sent Events (SSE) opcionales para streaming.
    Habilita la comunicación con servidores remotos y soporta métodos estándar
    de autenticación HTTP.
\end{itemize}

La capa de transporte abstrae los detalles de comunicación, permitiendo que el
mismo formato de mensajes JSON-RPC 2.0 funcione de manera idéntica sobre
cualquier mecanismo de transporte.

\subsubsection{Flujo de interacción}

Una interacción típica entre un cliente y un servidor MCP sigue una secuencia
bien definida \citep{mcpspec2024}:

\begin{enumerate}
    \item \textbf{Inicialización}: el cliente envía una solicitud
    \texttt{initialize} con su versión del protocolo y capacidades soportadas.
    El servidor responde con las suyas. Ambas partes negocian qué primitivas
    y funcionalidades estarán disponibles durante la sesión.

    \item \textbf{Descubrimiento de herramientas}: el cliente invoca
    \texttt{tools/list} para obtener la lista completa de herramientas
    disponibles, incluyendo nombre, descripción y schema de entrada
    (\texttt{inputSchema}) definido como JSON Schema para cada una.

    \item \textbf{Ejecución de herramientas}: cuando el LLM del host decide
    utilizar una herramienta, el cliente envía una solicitud
    \texttt{tools/call} con el nombre de la herramienta y sus argumentos. El
    servidor ejecuta la operación y retorna el resultado como un arreglo de
    objetos de contenido (texto, imágenes, recursos).

    \item \textbf{Actualización dinámica}: si la lista de herramientas cambia
    durante la sesión, el servidor envía una notificación
    \texttt{notifications/tools/list\_changed} y el cliente solicita la lista
    actualizada.
\end{enumerate}

\subsubsection{Adopción del ecosistema}

Desde su lanzamiento, MCP ha experimentado una adopción amplia tanto en la
industria como en la comunidad open-source. En el lanzamiento inicial,
Anthropic publicó la especificación abierta, SDKs para Python y TypeScript, y
servidores de referencia para Google Drive, Slack, GitHub, Git, Postgres y
Puppeteer \citep{anthropic2024mcp}. Empresas como Block y Apollo integraron
el protocolo en sus operaciones, mientras que plataformas de desarrollo como
Zed, Replit, Codeium y Sourcegraph lo adoptaron para mejorar las capacidades
contextuales de sus herramientas de IA.

A la fecha, MCP cuenta con soporte en aplicaciones como Claude Desktop,
ChatGPT, Visual Studio Code, Cursor y numerosos clientes adicionales
\citep{mcpspec2024}, consolidándose como el estándar de facto para la
integración de herramientas en sistemas basados en LLMs.
